\chapter{Conclusiones y líneas de trabajo futuro}

A lo largo del presente trabajo se ha desarrollado un análisis sistemático de diferentes metodologías de DV-QKD, contrastando sus fundamentos teóricos con las limitaciones físicas impuestas por los componentes de hardware comerciales. A continuación, se presentan las conclusiones globales derivadas de esta investigación, seguidas de una propuesta estructurada con las líneas de trabajo que definen la vanguardia de esta disciplina.

\section{Conclusiones}

La primera y más relevante conclusión de este estudio es que la seguridad teórica o incondicional de los sistemas cuánticos, basada en leyes de la física como el teorema de no clonación \citep{wootters1982}, se ve comprometida de forma crítica al trasladar los esquemas teóricos a entornos de implementación reales. La imperfección de las fuentes ópticas comerciales, dominadas por la emisión multi-fotónica estocástica de los pulsos coherentes atenuados, introduce vectores de vulnerabilidad severos, siendo el ataque de división del número de fotones la amenaza más destructiva en canales con altas pérdidas \citep{brassard2000pns, lutkenhaus2002}.

Frente a este escenario, la evaluación comparativa realizada en este Trabajo muestra que no existe un diseño único óptimo, sino que la selección de la arquitectura de red depende de un compromiso (\emph{trade-off}) de ingeniería entre la eficiencia del cribado clásico, la tolerancia a la atenuación del canal y la complejidad del hardware óptico utilizado:
\begin{itemize}
    \item Las arquitecturas basadas en \textbf{codificación estándar con intensidad constante} destacan por su simplicidad algorítmica y una alta eficiencia de cribado teórica del 50\%, pero su incapacidad intrínseca para segregar los pulsos monofotónicos de los multifotónicos limita drásticamente su distancia segura en redes reales, volviéndolas ineficaces ante tasas de pérdidas elevadas.
    \item Los esquemas basados en la \textbf{modulación dinámica de la intensidad media} de la fuente representan la solución más robusta y comercialmente viable para neutralizar los ataques de división de fotones. Al alterar de forma probabilística la potencia del láser, permiten calcular con rigor matemático cotas inferiores de ganancia monofotónica. Aunque incrementa la carga computacional en el postprocesamiento estadístico (\emph{finite-key analysis}), esta variante estira la distancia operativa de la fibra óptica a rangos metropolitanos avanzados con óptimas tasas de generación de clave (\emph{key rates}).
    \item Las variantes basadas en la \textbf{reconfiguración del anuncio clásico por conjuntos de estados} ofrecen una alternativa metodológica sumamente ingeniosa al modificar la etapa de postprocesamiento clásico para codificar la información en pares de estados en lugar de revelar las bases. Su gran ventaja es que aprovecha de manera segura los componentes multifotónicos sin requerir modificaciones en la potencia de emisión del hardware emisor. No obstante, su penalización en la eficiencia de cribado (reducida al 25\%) y su menor rendimiento en distancia frente a los métodos de modulación de intensidad las configuran como soluciones de nicho, ideales para enlaces de corta distancia con alta atenuación donde no se disponga de moduladores ópticos rápidos.
\end{itemize}

Finalmente, se concluye que los ataques de canal lateral (\emph{side-channel attacks}), tales como el cegado de detectores de avalancha mediante inyección de luz clásica \citep{lydersen2010}, constituyen el verdadero cuello de botella para la seguridad práctica de las redes de comunicaciones cuánticas contemporáneas. Esto demuestra que la validación de un sistema criptográfico cuántico no puede limitarse a la elegancia de sus demostraciones matemáticas, sino que requiere auditorías rigurosas de calibración y modelos de seguridad adaptados al comportamiento físico del hardware real.
